
%-------------------------------------------------%
% Pitch Markupp
%-------------------------------------------------%
\documentclass[aspectratio=169]{beamer}

\usepackage[english,brazil]{babel}
\usepackage[style=abnt,noslsn,justify]{biblatex}
\renewcommand*{\bibfont}{\small}
\usepackage{csquotes}
\addbibresource{referencias_trabalho.bib}

\usepackage{tikz}
\usepackage{pifont}
\usetikzlibrary{positioning, arrows.meta, shapes.geometric, shapes.symbols, decorations.pathmorphing}

\definecolor[named]{novaCor}{HTML}{303A52}
\definecolor{azul-seg}{RGB}{27, 85, 142}
\definecolor{verde-ok}{RGB}{34, 139, 87}
\definecolor{vermelho-x}{RGB}{180, 60, 60}
\definecolor{cinza-fundo}{RGB}{240, 240, 245}

\usetheme[
    foreground=azul-seg,
]{ifsclean}

% Comando para uma "tagline" em destaque no topo do slide
\newcommand{\tagline}[1]{%
    \begin{center}
        \colorbox{azul-seg!10}{%
            \begin{minipage}{0.95\linewidth}\centering
                \vspace{0.2em}
                {\bfseries\color{azul-seg} #1}
                \vspace{0.2em}
            \end{minipage}%
        }
    \end{center}
    \vspace{0.3em}
}

% -------------------------------------------------%
\title{Markupp}
\subtitle{Um hub para suas anotações e seus agentes de IA,\\com o dado sempre no seu servidor}
\author{
  Renato Freitas \and Nícolas Arthur \and Nicolas Pitz \and Gabriela Riedel
}
\institute{Instituto Federal de Santa Catarina}
\date{Maio de 2026}

% -------------------------------------------------%
\begin{document}

\begin{frame}[plain, noframenumbering]
    \maketitle
\end{frame}

\begin{frame}[plain, noframenumbering]{Licenciamento}
    \licenciamentoLivre
\end{frame}

\begin{frame}[plain, noframenumbering]{Sumário}
    \tableofcontents
\end{frame}

% -------------------------------------------------%
\section{Introdução}
% -------------------------------------------------%

\begin{frame}{O que é o Markupp?}
    \tagline{Um lugar único para suas anotações e para seus agentes de IA, instalado no seu próprio servidor.}

    \begin{columns}[T]
        \begin{column}{0.58\textwidth}
            \begin{itemize}\small
                \item \textbf{Servidor próprio}: você instala no seu computador ou na sua nuvem. Sem mensalidade, sem dependência de terceiros.
                \item \textbf{Formato aberto (Markdown)}: arquivos de texto simples, que você pode ler em qualquer editor.
                \item \textbf{Interface universal (API)}: o mesmo conjunto de notas é acessado pelo \textbf{usuário} e por \textbf{qualquer agente de IA}.
                \item \textbf{Open-source}: código aberto e auditável, sem ficar refém de uma empresa.
            \end{itemize}
        \end{column}
        \begin{column}{0.40\textwidth}
            \centering
            \begin{tikzpicture}[scale=0.80, every node/.style={transform shape},
                node distance=0.5cm and 0.3cm,
                box/.style={draw=azul-seg, thick, rounded corners, minimum width=1.8cm, minimum height=0.6cm, align=center, font=\scriptsize},
                hub/.style={draw=azul-seg, very thick, fill=azul-seg!15, rounded corners, minimum width=2.2cm, minimum height=0.9cm, align=center, font=\small\bfseries},
                arrow/.style={-Stealth, thick, azul-seg}
            ]
                \node[hub] (hub) {Markupp\\\scriptsize (seu servidor)};
                \node[box, above=of hub] (user) {Você};
                \node[box, below left=of hub] (obs) {Obsidian};
                \node[box, below right=of hub] (ai) {Agente IA};
                \draw[arrow, <->] (user) -- (hub);
                \draw[arrow, <->] (obs) -- (hub);
                \draw[arrow, <->] (ai) -- (hub);
            \end{tikzpicture}

            \vspace{0.3em}
            \tiny\itshape Uma base, vários clientes.
        \end{column}
    \end{columns}
\end{frame}

% -------------------------------------------------%
\section{Tamanho de mercado}
% -------------------------------------------------%

\begin{frame}{Um mercado em forte expansão}
    \tagline{O Markupp cruza dois mercados que somam US\$ 24 bi em 2026 e crescem mais de 20\% ao ano.}

    \begin{columns}[T]
        \begin{column}{0.50\textwidth}
            \begin{itemize}\footnotesize
                \item \textbf{Apps de anotações}: US\$ 13{,}3 bi em 2026 \cite{tbrc_notetaking}.
                \item \textbf{Agentes de IA}: US\$ 10{,}9 bi, \textasciitilde 46\% a.a. \cite{grandview_aiagents}.
                \item Gartner: \textbf{40\%} dos apps empresariais terão agentes de IA até o fim de 2026 \cite{gartner_aiagents}.
            \end{itemize}

            \vspace{0.3em}
            \textbf{Por que open-source agora?}
            \begin{itemize}\footnotesize
                \item LGPD/GDPR pressionam empresas a controlar seus dados \cite{dreamhost_selfhosted}.
                \item Docker tornou instalar servidores próprios trivial.
            \end{itemize}
        \end{column}
        \begin{column}{0.46\textwidth}
            \centering\textbf{Anotações + agentes de IA (US\$ bi)}\\[0.3em]
            \begin{tikzpicture}[scale=0.85]
                \draw[->, thick] (0,0) -- (5.6,0) node[right, font=\tiny] {ano};
                \draw[->, thick] (0,0) -- (0,3.0) node[above, font=\tiny] {US\$ bi};
                \draw[azul-seg, very thick] (0.3,0.5) -- (1.6,1.0) -- (2.7,1.5) -- (3.8,2.1) -- (4.9,2.6);
                \foreach \x/\y/\v in {0.3/0.5/19, 1.6/1.0/24, 2.7/1.5/30, 3.8/2.1/38, 4.9/2.6/48} {
                    \fill[azul-seg] (\x,\y) circle (1.5pt);
                    \node[font=\scriptsize, azul-seg, above=1pt] at (\x,\y) {\v};
                }
                \foreach \x/\ano in {0.3/25, 1.6/26, 2.7/27, 3.8/28, 4.9/29} {
                    \node[font=\scriptsize, anchor=north] at (\x,0) {\ano};
                }
            \end{tikzpicture}

            \vspace{0.3em}
            \footnotesize\itshape Por sermos open-source: medimos \textbf{potencial de adoção}, não receita.
        \end{column}
    \end{columns}
\end{frame}

% -------------------------------------------------%
\section{Problema}
% -------------------------------------------------%

\begin{frame}{O problema: conhecimento fragmentado}
    \tagline{Suas anotações estão espalhadas em várias ferramentas; nenhuma delas conversa com seu agente de IA.}

    \begin{columns}[T]
        \begin{column}{0.50\textwidth}
            \textbf{Cenário típico}
            \begin{itemize}\footnotesize
                \item Você anota uma ideia no \textbf{Notion}, faz \emph{brainstorm} no \textbf{ChatGPT} e cola um link no \textbf{Slack}.
                \item No dia seguinte, nem você nem a IA conseguem reunir tudo.
                \item SaaS retêm dados na nuvem; ferramentas locais (Obsidian) protegem privacidade, mas \textbf{não expõem interface} para outros programas.
            \end{itemize}

            \vspace{0.3em}
            \footnotesize\emph{Resultado}: retrabalho, perda de contexto e zero simbiose com a IA.
        \end{column}
        \begin{column}{0.48\textwidth}
            \centering
            \begin{tikzpicture}[scale=0.78, every node/.style={transform shape},
                tool/.style={draw=vermelho-x!70, thick, rounded corners, fill=vermelho-x!5,
                             minimum width=1.3cm, minimum height=0.55cm, align=center, font=\tiny},
                user/.style={draw=azul-seg, very thick, circle, fill=azul-seg!15,
                             minimum size=1.0cm, align=center, font=\scriptsize\bfseries},
                wave/.style={->, decorate, decoration={snake, amplitude=0.5mm, segment length=2mm,
                             post length=1.0mm}, vermelho-x!80, thick}
            ]
                \node[user] (u) at (0,0) {Você};
                \node[tool] (n)  at (-2.6, 1.3) {Notion};
                \node[tool] (gpt) at (-1.3, 1.8) {ChatGPT};
                \node[tool] (sl)  at ( 0.2, 2.0) {Slack};
                \node[tool] (gd)  at ( 1.7, 1.8) {GDocs};
                \node[tool] (ap)  at ( 2.8, 1.3) {Apple\\Notes};
                \node[tool] (em)  at ( 2.8,-1.1) {e-mail};
                \node[tool] (ob)  at ( 1.3,-1.9) {Obsidian};
                \node[tool] (ag)  at (-1.3,-1.9) {Agente?};
                \node[tool] (dr)  at (-2.8,-1.1) {Drive};
                \foreach \t in {n,gpt,sl,gd,ap,em,ob,ag,dr}
                    \draw[wave] (u) -- (\t);
                \node[font=\tiny\itshape, vermelho-x] at (0,-2.6) {silos isolados, sem visão unificada};
            \end{tikzpicture}
        \end{column}
    \end{columns}
\end{frame}

% -------------------------------------------------%
\section{Solução}
% -------------------------------------------------%

\begin{frame}{A solução: Markupp}
    \tagline{Um servidor open-source que centraliza seus arquivos em Markdown e os expõe por uma interface aberta.}

    \begin{columns}[T]
        \begin{column}{0.55\textwidth}
            \textbf{Em uma frase:} um back-end único guarda suas notas; vários clientes conversam com ele pela mesma interface aberta.
            \begin{itemize}\small
                \item \textbf{Criar, ler, editar e excluir} notas pelo plugin do Obsidian.
                \item \textbf{Interface aberta (REST)}: qualquer agente de IA lê e escreve nas mesmas notas.
                \item \textbf{No seu servidor}, sem nuvem terceirizada. Código auditável.
            \end{itemize}
        \end{column}
        \begin{column}{0.42\textwidth}
            \textbf{Já validado no MVP}
            \begin{itemize}\footnotesize
                \item[\textcolor{verde-ok}{\ding{51}}] Back-end (Go) com armazenamento
                \item[\textcolor{verde-ok}{\ding{51}}] Criar, ler, editar e excluir notas
                \item[\textcolor{verde-ok}{\ding{51}}] Plugin Obsidian: criação
                \item[\textcolor{verde-ok}{\ding{51}}] Testes + integração contínua
                \item[$\circ$] Plugin: edição/exclusão/rename (em curso)
            \end{itemize}

            \vspace{0.3em}
            \footnotesize \emph{Stack:} Go, PostgreSQL, Docker, TypeScript.
        \end{column}
    \end{columns}
\end{frame}

% -------------------------------------------------%
\section{Concorrência}
% -------------------------------------------------%

\begin{frame}{Como nos posicionamos}
    \tagline{Cada concorrente cobre parte do problema. O Markupp é o único que combina os 4 atributos.}

    \scriptsize
    \begin{center}
    \begin{tabular}{|l|c|c|c|c|}
        \hline
        \textbf{Solução} & \textbf{No seu servidor} & \textbf{Hub central + API} & \textbf{Markdown} & \textbf{Pensado p/ IA} \\
        \hline
        Notion       & \textcolor{vermelho-x}{\ding{55}} (SaaS) & \textcolor{verde-ok}{\ding{51}} (proprietária) & \textcolor{vermelho-x}{\ding{55}} & \textcolor{vermelho-x}{\ding{55}} \\
        \hline
        Obsidian     & \textcolor{verde-ok}{\ding{51}} (local) & \textcolor{vermelho-x}{\ding{55}} & \textcolor{verde-ok}{\ding{51}} & \textcolor{vermelho-x}{\ding{55}} \\
        \hline
        Logseq       & \textcolor{verde-ok}{\ding{51}} (local) & \textcolor{vermelho-x}{\ding{55}} & \textcolor{verde-ok}{\ding{51}} & \textcolor{vermelho-x}{\ding{55}} \\
        \hline
        AnyType      & \textcolor{verde-ok}{\ding{51}} & \textcolor{vermelho-x}{\ding{55}} (formato próprio) & \textcolor{vermelho-x}{\ding{55}} & \textcolor{vermelho-x}{\ding{55}} \\
        \hline
        \textbf{Markupp} & \textcolor{verde-ok}{\ding{51}} & \textcolor{verde-ok}{\ding{51}} (REST aberta) & \textcolor{verde-ok}{\ding{51}} & \textcolor{verde-ok}{\ding{51}} \\
        \hline
    \end{tabular}
    \end{center}

    \vspace{0.4em}
    \normalsize
    \begin{itemize}\footnotesize
        \item \textbf{Notion}: centraliza, mas é SaaS fechado; seu dado fica na nuvem deles.
        \item \textbf{Obsidian/Logseq}: controle local, mas não compartilham a base com agentes.
        \item \textbf{Markupp}: hub central + controle do dado + abertura para a IA.
    \end{itemize}
\end{frame}

% -------------------------------------------------%
\section{Time}
% -------------------------------------------------%

\begin{frame}{Time}
    \tagline{Quatro desenvolvedores, cada um responsável por um papel de processo.}

    \begin{columns}[T]
        \begin{column}{0.5\textwidth}
            \begin{itemize}\small
                \item \textbf{Renato Freitas}\\Desenvolvedor + \emph{Arquiteto de Software}: desenho da API e do back-end em Go.
                \vspace{0.3em}
                \item \textbf{Nícolas Arthur}\\Desenvolvedor + \emph{DevOps / Infra}: contêineres, integração contínua e infraestrutura self-hosted.
            \end{itemize}
        \end{column}
        \begin{column}{0.5\textwidth}
            \begin{itemize}\small
                \item \textbf{Nicolas Pitz}\\Desenvolvedor + \emph{Qualidade}: testes ($\geq 80\%$ de cobertura) e \emph{code review}.
                \vspace{0.3em}
                \item \textbf{Gabriela Riedel}\\Desenvolvedora + \emph{Scrum Master}: planejamento, métricas e cadência das sprints.
            \end{itemize}
        \end{column}
    \end{columns}

    \vspace{0.6em}
\end{frame}

% -------------------------------------------------%
\section{Próximos passos}
% -------------------------------------------------%

\begin{frame}{Próximos passos}
    \tagline{Um \emph{roadmap} em três horizontes: terminar o MVP, refinar o produto, virar plataforma de IA.}

    \begin{columns}[T]
        \begin{column}{0.32\textwidth}
            \textbf{Curto prazo}\\\footnotesize \emph{fechar o MVP}
            \begin{itemize}\footnotesize
                \item Plugin Obsidian completo: edição, exclusão e renomear.
                \item Manter cobertura $\geq 80\%$.
            \end{itemize}
        \end{column}
        \begin{column}{0.32\textwidth}
            \textbf{Médio prazo}\\\footnotesize \emph{refinar o produto}
            \begin{itemize}\footnotesize
                \item Pastas e organização hierárquica.
                \item Busca por título e conteúdo.
            \end{itemize}
        \end{column}
        \begin{column}{0.32\textwidth}
            \textbf{Longo prazo}\\\footnotesize \emph{diferencial competitivo}
            \begin{itemize}\footnotesize
                \item Busca semântica (vetorial).
                \item Conexão com agentes via MCP (protocolo padrão para IA).
                \item Comunidade de contribuidores.
            \end{itemize}
        \end{column}
    \end{columns}

    \vspace{0.8em}
    \begin{center}
        \large\bfseries\color{azul-seg}
        Um segundo cérebro aberto, soberano e pronto para a era dos agentes.
    \end{center}
\end{frame}

% -------------------------------------------------%
\begin{frame}[allowframebreaks, noframenumbering]{Referências}
    \nocite{*}
    \printbibliography
\end{frame}

\end{document}
